\section{Transmission lines: the distributed circuit}
\label{sec:tline}

\noindent\emph{Status: settled material, included because everything after it depends on it. Nothing in this
section is a claim about the solver. What is derived here -- the propagation constant, the characteristic
impedance, the reflection coefficient, the half-wavelength ripple, the quarter-wave identity, and the
operational meaning of the effective permittivity -- is used without further argument by
Section~\ref{sec:microstrip}, and the three chapters that follow this one exist so that the measurement
language of the later sections is derived rather than assumed.}

The reader arriving from Section~\ref{sec:impl} has a solver that computes fields, and a box that supports a
mode which travels along a strip. Neither of those is yet a measurement. A measurement needs a vocabulary in
which the field inside the box can be summarised as a handful of complex numbers that a network analyser would
also have produced, on a bench, from the same physical object. That vocabulary is built in three steps. This
section builds the first: the description of a guiding structure by two conductors and four numbers per metre,
from which a propagation constant and a characteristic impedance follow. Section~\ref{sec:sparams} builds the
second: the description of a whole network by waves entering and leaving its ports. Section~\ref{sec:smith}
builds the third: the graphical representation of one complex number, the reflection coefficient, in which the
arithmetic of matching becomes geometry.

The order matters, and it is the reason these three sections stand between the implementation and the
microstrip. A scattering parameter is defined as a ratio of wave amplitudes at a port, and a wave at a port is
a travelling-wave solution of the equations derived below, evaluated at a reference impedance defined below.
Quoting an $S_{11}$ for a structure whose guided wavelength and whose characteristic impedance have not been
established is quoting a number with two free definitions inside it. The microstrip sections that follow are
unusually careful about exactly this point, and they can afford to be only because the definitions are fixed
here.

\subsection{When a wire stops being a node}
\label{sec:tline-lumped}

Lumped-circuit analysis rests on an assumption so familiar that it is rarely stated: that the voltage at every
point of a connection is the same, so that a wire may be replaced by a mathematical node. The assumption is
false at every frequency, and it becomes a useful approximation below a size that depends on the guided
wavelength. The size of the failure is a phase difference. If a structure of length $l$ carries a wave of
guided wavelength $\lambda_g$, the phase of the field at one end relative to the other is
\begin{equation}
  \Delta\varphi = 2\pi\,\frac{l}{\lambda_g}\ \text{radians},
  \label{eq:tline-dphi}
\end{equation}
and a lumped model asserts $\Delta\varphi = 0$. Nothing in circuit theory fails gradually: the assertion is
simply wrong by $\Delta\varphi$, and whether that is acceptable is a statement about the accuracy wanted from
the model, not about the circuit.

What makes the criterion confusing is that it is usually stated as a fraction of a wavelength rather than as a
phase. The common rule of thumb is that a structure may be treated as lumped when $l < \lambda_g/10$, which by
\eqref{eq:tline-dphi} admits a phase error of $36^\circ$; a rule of $l < \lambda_g/20$ admits $18^\circ$; an
accuracy target of one per cent in a resonance frequency corresponds to a considerably shorter structure.
Table~\ref{tab:tline-lumped} tabulates the lengths at which each of those conventions fails for a line on a
substrate of $\eps_{\mathrm{eff}} = 3.3$, which is the value the reference box of
Section~\ref{sec:msres-box} is close to. All of them are small at microwave frequencies, and all of them are
large compared with a printed trace. That coincidence is the whole reason printed circuits at microwave
frequencies are transmission-line problems.

\begin{table}[H]
  \centering
  \caption{Lengths at which lumped-circuit analysis stops being defensible, for a guided wavelength on a
  substrate of $\eps_{\mathrm{eff}} = 3.3$ ($\sqrt{\eps_{\mathrm{eff}}} = 1.8166$, so
  $\lambda_g = \lambda_0/1.8166$). The right-hand column is the length at which \eqref{eq:tline-dphi} reaches
  $36^\circ$, the phase error hidden by the folk rule $l<\lambda_g/10$. A 10 mm trace is already a distributed
  structure by the middle of the table.}
  \label{tab:tline-lumped}
  \small
  \begin{tabular}{@{}rrrrr@{}}
    \toprule
    $f$ & $\lambda_0$ & $\lambda_g$ & $\lambda_g/10$ & $\lambda_g/63$ ($1^\circ$ per cell) \\
    \midrule
    $1$ MHz     & $300$ m      & $165$ m      & $16.5$ m     & $2.62$ m \\
    $100$ MHz   & $3.00$ m     & $1.65$ m     & $165$ mm     & $26.2$ mm \\
    $1$ GHz     & $300$ mm     & $165$ mm     & $16.5$ mm    & $2.62$ mm \\
    $10$ GHz    & $30.0$ mm    & $16.5$ mm    & $1.65$ mm    & $262\ \mu$m \\
    \bottomrule
  \end{tabular}
\end{table}

The criterion is sometimes restated in the time domain, where it reads more directly and is more useful for
digital work. A signal with rise time $t_r$ has content up to a bandwidth of order $1/t_r$, so the shortest
wavelength present is of order $v_p t_r$, and the lumped assumption holds while the structure is short against
that: $l \ll v_p t_r$, with $l < v_p t_r/6$ the usual working form. On a printed board with
$v_p = 2.0\times10^{8}$ m/s, a $100$ ps edge is $20$ mm of propagation, so traces longer than a few millimetres
already carry a visible phase gradient. The two forms of the criterion are the same statement: what matters is
the electrical length in wavelengths, and neither the physical size of the object nor the frequency alone
determines it.

\subsection{The infinitesimal section and the telegrapher's equations}
\label{sec:tline-telegrapher}

The distributed description replaces the wire by a limit of short sections, each of which is a lumped circuit
small enough for the previous subsection's criterion to hold locally. Consider a length $\dd z$ of a
two-conductor line, with the two conductors carrying currents $I(z,t)$ and $-I(z,t)$ and a voltage $V(z,t)$
between them. The section is modelled by four quantities, all referred to unit length: a series resistance
$R'$ and a series inductance $L'$ in the conductor path, and a shunt conductance $G'$ and a shunt capacitance
$C'$ between the conductors. Their units are ohm per metre, henry per metre, siemens per metre and farad per
metre, and the primes are there to distinguish them from circuit elements measured in ohms, henrys, siemens and
farads. A line of length $L$ with no reflections has a total series inductance $L' L$, and that is all the
prime means.

Kirchhoff's voltage law around the loop formed by the two conductors over the section, divided by $\dd z$ and
taken to the limit, gives the loss of voltage along the line as the drop across the series elements; the
current law at the node between the section and the next gives the loss of current along the line as the
current diverted through the shunt elements. The two statements are
\begin{equation}
  \frac{\partial V}{\partial z} = -\left(R' I + L' \frac{\partial I}{\partial t}\right),
  \qquad
  \frac{\partial I}{\partial z} = -\left(G' V + C' \frac{\partial V}{\partial t}\right),
  \label{eq:tline-telegrapher}
\end{equation}
with the sign convention that $z$ increases in the direction of the incident wave and that the voltages and
currents are of the forward-travelling wave. The first equation is a statement about flux: the voltage falls
along the line at a rate set by the resistive drop and by the rate of change of the magnetic flux threading
the section. The second is a statement about charge: the current falls at a rate set by the leakage through
the dielectric and by the rate of change of the charge stored on the section. Nothing in
\eqref{eq:tline-telegrapher} refers to a field, a grid, or a mode; it is an exact consequence of the
definitions of $V$ and $I$ once those are taken as line integrals of $\mathbf{E}$ and $\mathbf{H}$ over a
transverse plane, and it is exact only to the extent that the transverse field pattern does not change with
$z$.

For harmonic excitation the time dependence is $e^{j\omega t}$ and the two partial differential equations
collapse to a pair of ordinary ones,
\begin{equation}
  \frac{\dd V}{\dd z} = -(R' + j\omega L')\,I, \qquad
  \frac{\dd I}{\dd z} = -(G' + j\omega C')\,V .
  \label{eq:tline-harmonic}
\end{equation}
Differentiating either and substituting the other decouples them into a second-order equation for each
quantity,
\begin{equation}
  \frac{\dd^2 V}{\dd z^2} = \gamma^2 V, \qquad
  \gamma = \sqrt{(R' + j\omega L')(G' + j\omega C')} = \alpha + j\beta,
  \label{eq:tline-propgamma}
\end{equation}
with $\gamma$ the propagation constant in reciprocal metres, $\alpha$ the attenuation constant in nepers per
metre and $\beta$ the phase constant in radians per metre. The sign convention is chosen so that a forward wave
is $e^{-\gamma z}$: a solution that decays as it travels has $\alpha > 0$, and a solution that accumulates phase
has $\beta > 0$. This is the same convention as Section~\ref{sec:microstrip} uses for the mode of a strip,
\eqref{eq:ms-hybrid}, and the same $\alpha$ and $\beta$ that the station fits of
Section~\ref{sec:ms-estimators} return, so the two descriptions can be compared without a change of variable.

The general solution is a sum of two exponentials,
\begin{equation}
  V(z) = V^{+}e^{-\gamma z} + V^{-}e^{+\gamma z}, \qquad
  I(z) = \frac{V^{+}}{Z_0}e^{-\gamma z} - \frac{V^{-}}{Z_0}e^{+\gamma z},
  \label{eq:tline-solution}
\end{equation}
with $Z_0$ fixed by substituting \eqref{eq:tline-solution} into \eqref{eq:tline-harmonic} and requiring both
signs to be consistent,
\begin{equation}
  Z_0 = \sqrt{\frac{R' + j\omega L'}{G' + j\omega C'}} .
  \label{eq:tline-z0-general}
\end{equation}
The minus sign in front of the backward current is not a convention that may be chosen: it follows from the
fact that the backward wave carries power in the $-z$ direction, so its current must be negative with respect
to the same reference direction as the forward wave. That sign is the origin of every subsequent statement
about reflection, and getting it wrong is the most common way a derivation of $S_{11}$ ends up with the
reciprocal of the right answer.

\subsection{The lossless line, and where $Z_0$ and $v_p$ come from}
\label{sec:tline-lossless}

Setting $R' = G' = 0$ in \eqref{eq:tline-propgamma} and \eqref{eq:tline-z0-general} gives the lossless line, on
which every quantity of interest becomes a real square root and every formula can be checked by hand,
\begin{equation}
  \gamma = j\omega\sqrt{L'C'} = j\beta, \qquad
  \alpha = 0, \qquad
  \beta = \omega\sqrt{L'C'}, \qquad
  Z_0 = \sqrt{\frac{L'}{C'}} .
  \label{eq:tline-lossless}
\end{equation}
Three consequences deserve to be spelled out, because the rest of the report uses all three. The first is that
$Z_0$ is real and frequency-independent on a lossless line: it is a property of the geometry and the materials
through $L'$ and $C'$, not of the drive. The second is the phase velocity,
\begin{equation}
  v_p = \frac{\omega}{\beta} = \frac{1}{\sqrt{L'C'}},
  \label{eq:tline-vp}
\end{equation}
which is likewise real, so that a lossless line delays a signal without distorting it. The third is that the
two quantities are not independent: \eqref{eq:tline-lossless} and \eqref{eq:tline-vp} can be solved for the
per-unit-length parameters in terms of the two measurable ones,
\begin{equation}
  L' = \frac{Z_0}{v_p}, \qquad C' = \frac{1}{Z_0 v_p}, \qquad L'C' = \frac{1}{v_p^{2}}, \qquad
  \frac{L'}{C'} = Z_0^{2} .
  \label{eq:tline-lc}
\end{equation}
A $50\ \Ohm$ line on a substrate that slows the wave to $2.0\times10^{8}$ m/s therefore has
$L' = 250$ nH/m and $C' = 100$ pF/m, and a reader who multiplies those two numbers gets
$2.5\times10^{-17}\,\mathrm{s}^2/\mathrm{m}^2$, which is the reciprocal square of the velocity. The check is
worth performing once: it fixes the orders of magnitude that the rest of the section quotes, and it is the same
arithmetic by which a measured $\beta$ is converted into an effective permittivity below.

The derivation of $Z_0$ is worth doing explicitly rather than quoting \eqref{eq:tline-z0-general}, because it
shows what $Z_0$ is a ratio \emph{of}. Take the lossless first equation of \eqref{eq:tline-harmonic} with the
forward wave alone, $V = V^{+}e^{-j\beta z}$. Then $\dd V/\dd z = -j\beta V^{+}e^{-j\beta z} = -j\omega L' I$,
so
\begin{equation}
  I = \frac{\beta}{\omega L'}\,V = \frac{\omega\sqrt{L'C'}}{\omega L'}\,V
    = \sqrt{\frac{C'}{L'}}\,V = \frac{V}{Z_0},
  \label{eq:tline-z0-derivation}
\end{equation}
using $\beta = \omega\sqrt{L'C'}$ from \eqref{eq:tline-lossless}. The characteristic impedance is therefore the
ratio of the voltage to the current \emph{of a single travelling wave}, and nothing else. It is not the ratio
that would be measured at a plane where two waves are present; that ratio is the input impedance, which is a
property of the termination as much as of the line and is derived in Section~\ref{sec:tline-zin}. Confusing the
two is the first of the three confusions catalogued in Section~\ref{sec:ms-z0}, where the same distinction is
made for a microstrip and where the spread between the definitions is used as a measurement rather than
dismissed as an annoyance.

A lossy line is not qualitatively different, and the low-loss limit is worth recording because the microstrip
sections refer to attenuation in decibels per millimetre and the conversion has to come from somewhere. For
$R' \ll \omega L'$ and $G' \ll \omega C'$, expanding \eqref{eq:tline-propgamma} and
\eqref{eq:tline-z0-general} to first order gives
\begin{equation}
  \alpha \simeq \tfrac12\left(\frac{R'}{Z_0} + G' Z_0\right), \qquad
  \beta \simeq \omega\sqrt{L'C'}, \qquad
  Z_0 \simeq \sqrt{\frac{L'}{C'}}\left[1 + \frac{j}{2}\left(\frac{G'}{\omega C'}
  - \frac{R'}{\omega L'}\right)\right],
  \label{eq:tline-lowloss}
\end{equation}
so the attenuation has two additive mechanisms, one in the conductor and one in the dielectric, and the
characteristic impedance acquires a small imaginary part. The neper is the natural unit of $\alpha$; one neper
per metre is $8.6859$ dB per metre, or $8.6859\times10^{-3}$ dB per millimetre, which is the conversion used
wherever an attenuation appears in this report in decibels.

\subsection{The reflection coefficient}
\label{sec:tline-gamma}

A line that is not terminated in $Z_0$ has nowhere to send part of the wave that arrives at its end, and that
part comes back. This is not a metaphor: it is the consequence of the two boundary conditions at the load, and it
is worth deriving once so that the symbol $\Gamma$ never again appears as an incantation. Write the general
solution of the lossless line as a sum of a forward wave travelling in $+z$ and a backward wave travelling in
$-z$,
\begin{equation}
  V(z) = V^{+} e^{-j\beta z} + V^{-} e^{j\beta z}, \qquad
  I(z) = \frac{V^{+}}{Z_0} e^{-j\beta z} - \frac{V^{-}}{Z_0} e^{j\beta z},
  \label{eq:tline-two-waves}
\end{equation}
where the sign of the second current term is the only thing that distinguishes a wave carrying power towards
$+z$ from one carrying power towards $-z$: for the backward wave the current reverses with respect to the
voltage. Let the load sit at $z=0$, so that the ratio $V(0)/I(0)$ must equal the load impedance $Z_L$:
\begin{equation}
  Z_L = \frac{V^{+} + V^{-}}{V^{+}/Z_0 - V^{-}/Z_0}
      = Z_0\,\frac{1 + \Gamma_L}{1 - \Gamma_L},
  \qquad \Gamma_L \equiv \frac{V^{-}}{V^{+}} .
  \label{eq:tline-gamma-def}
\end{equation}
Inverting \eqref{eq:tline-gamma-def} is two lines of algebra and gives the reflection coefficient at the load,
\begin{equation}
  \Gamma_L = \frac{Z_L - Z_0}{Z_L + Z_0},
  \label{eq:tline-gamma}
\end{equation}
which is the quantity the rest of this report uses constantly. Three cases fix its meaning. A load equal to
$Z_0$ gives $\Gamma_L=0$: the line is \emph{matched}, no backward wave exists, and the voltage magnitude is the
same at every station. An open circuit gives $\Gamma_L=+1$: the backward wave equals the forward wave in both
amplitude and phase, which at the open end doubles the voltage and cancels the current, as an open circuit must.
A short circuit gives $\Gamma_L=-1$: the backward wave is inverted, which zeroes the voltage and doubles the
current. Between these, $|\Gamma_L|$ measures how badly the load is matched.

There is a property of \eqref{eq:tline-gamma} that matters more than any of the special cases: for any passive
load, $|\Gamma_L|\le 1$. The proof is a statement about the right half-plane. The map $Z_L\mapsto\Gamma_L$ sends
$\Re\{Z_L\}\ge 0$ into the closed unit disc, because
\begin{equation}
  |\Gamma_L|^2 = \frac{(R_L-Z_0)^2 + X_L^2}{(R_L+Z_0)^2 + X_L^2} \le 1
  \iff (R_L-Z_0)^2 \le (R_L+Z_0)^2 \iff 0 \le 4 R_L Z_0 ,
  \label{eq:tline-gamma-passive}
\end{equation}
which holds exactly when the resistive part $R_L$ is non-negative, i.e. when the load can only dissipate and
never generate. Equality holds only for $R_L=0$, the purely reactive load. An active load with $R_L<0$ produces
$|\Gamma_L|>1$ and therefore a backward wave larger than the forward one, which is the distributed-language
statement that the load is a source. The inequality \eqref{eq:tline-gamma-passive} is the same statement that
will be read geometrically in Section~\ref{sec:smith} as ``the passive half-plane lands inside the unit circle'',
and it is the same statement that reappears in Section~\ref{sec:sparams} as the passivity condition on a
scattering matrix.

\subsection{The standing wave, and why the ripple has period $\lambda_g/2$}
\label{sec:tline-ripple}

The sum of a forward and a backward wave of the same frequency is a standing wave, and the point of this
subsection is to derive -- not assert -- the distance over which it repeats. Let the reflection coefficient at
$z=0$ be $\Gamma_L=|\Gamma|e^{j\theta}$. At a distance $z$ towards the generator the backward wave has travelled
an extra $2z$ relative to the forward one: $z$ going away from the load and another $z$ coming back. One
factor of $e^{-j\beta z}$ per leg, so
\begin{equation}
  \Gamma(z) = \Gamma_L e^{-2j\beta z} = |\Gamma| e^{j(\theta - 2\beta z)} ,
  \label{eq:tline-gamma-z}
\end{equation}
and \eqref{eq:tline-gamma-z} already contains the answer, since the phase of $\Gamma$ advances by $2\beta$ per
metre, not $\beta$. Now compute the voltage magnitude, which is what a probe measures:
\begin{align}
  |V(z)|^2 &= \left(V^{+} e^{-j\beta z} + V^{-} e^{+j\beta z}\right)
              \left(V^{+*} e^{+j\beta z} + V^{-*} e^{-j\beta z}\right) \nonumber\\
           &= |V^{+}|^2 + |V^{-}|^2 + V^{+}V^{-*} e^{-2j\beta z} + V^{+*}V^{-} e^{+2j\beta z} \nonumber\\
           &= |V^{+}|^2\left[1 + |\Gamma|^2 + 2|\Gamma|\cos\!\left(2\beta z - \theta\right)\right].
  \label{eq:tline-ripple}
\end{align}
The derivation is finished: in \eqref{eq:tline-ripple} the only dependence on position is inside the cosine,
and the argument is $2\beta z-\theta$. A cosine returns to its value when its argument advances by $2\pi$, so
the pattern repeats over a distance $\Delta z$ satisfying $2\beta\,\Delta z = 2\pi$, that is
\begin{equation}
  \Delta z = \frac{\pi}{\beta} = \frac{\pi}{2\pi/\lambda_g} = \frac{\lambda_g}{2} .
  \label{eq:tline-ripple-period}
\end{equation}
The ripple period on a transmission line is half a guided wavelength, and the factor of two in the argument is
the whole reason. Physically it is the round trip: to see the same relative phase between the two travelling
waves, the probe must move by half a wavelength, because moving by a full wavelength advances the forward wave
by $2\pi$ and the backward wave by $-2\pi$, a relative change of $4\pi$ -- two full cycles of the interference
term. The same factor of two governs the reference-plane rotation of Section~\ref{sec:sp-refplane} and the
rotation of the Smith chart in Section~\ref{sec:smith-rotation}; all three are the same equation wearing
different hats.

Two landmarks follow immediately. Maxima of $|V|$ occur where $\cos(2\beta z-\theta)=+1$, i.e. at
$2\beta z = \theta + 2n\pi$, which are spaced by $\Delta z=\pi/\beta=\lambda_g/2$; minima sit halfway between
them and are also spaced by $\lambda_g/2$. Their ratio is the voltage standing wave ratio,
\begin{equation}
  \mathrm{VSWR} = \frac{|V|_{\max}}{|V|_{\min}}
               = \frac{1+|\Gamma|}{1-|\Gamma|} ,
  \qquad |\Gamma| = \frac{\mathrm{VSWR}-1}{\mathrm{VSWR}+1},
  \label{eq:tline-vswr}
\end{equation}
obtained by setting the cosine to $+1$ and $-1$ in \eqref{eq:tline-ripple} and taking the square roots. The
VSWR is therefore nothing but a monotone relabelling of $|\Gamma|$: it carries exactly as much information and
is quoted because it can be read off a slotted line with a detector diode that responds to magnitude only. And
the ``envelope'' of the standing wave, the pair of curves $\pm 20\log_{10}(1\pm|\Gamma|)$ in decibels around the
matched level, is what the estimator comparison of Section~\ref{sec:ms-estimators} has to fight. A concrete
number makes the difficulty explicit. With $|\Gamma|=0.0060$ -- a return loss of 44.4 dB, which sounds like an
excellent match -- \eqref{eq:tline-vswr} gives a VSWR of $1.0121$ and an envelope of $\pm1.21$ per cent. That is
precisely the ``$\pm1.2$ per cent'' figure quoted in Section~\ref{sec:ms-criterion}, and it is worth
understanding that this is not an error but the exact, unavoidable amplitude modulation that any backward wave
imposes on the standing wave. An estimator that assumes a constant envelope is wrong by 1.2 per cent before it
begins; an estimator that knows the period is $\lambda_g/2$ can average the ripple out, which is the entire
content of the estimators compared later.

\subsection{Input impedance, and the quarter- and half-wave identities}
\label{sec:tline-zin}

The impedance seen looking into a length $d$ of line terminated in $Z_L$ is the ratio $V(-d)/I(-d)$, and it is
the workhorse of everything that follows. Substituting \eqref{eq:tline-gamma-z} into
$Z(d)=Z_0(1+\Gamma(d))/(1-\Gamma(d))$ and dividing through by $e^{-j\beta d}$ gives
\begin{equation}
  Z_{\mathrm{in}}(d) = Z_0\,\frac{Z_L + j Z_0 \tan\beta d}{Z_0 + j Z_L \tan\beta d} ,
  \label{eq:tline-zin}
\end{equation}
a M\"obius transformation of $Z_L$ whose only frequency dependence is through $\tan\beta d$. Two values of that
tangent deserve their own names. At $\beta d=\pi/2$, i.e. $d=\lambda_g/4$, the tangent is infinite and
\eqref{eq:tline-zin} reduces to
\begin{equation}
  Z_{\mathrm{in}}(\lambda_g/4) = \frac{Z_0^2}{Z_L} ,
  \label{eq:tline-quarter}
\end{equation}
the quarter-wave transformer: a short circuit becomes an open circuit and an open circuit becomes a short, and
a real load $R_L$ is transformed into the real impedance $Z_0^2/R_L$. At $\beta d=\pi$, i.e. $d=\lambda_g/2$,
the tangent vanishes and
\begin{equation}
  Z_{\mathrm{in}}(\lambda_g/2) = Z_L ,
  \label{eq:tline-half}
\end{equation}
so a half-wavelength of line is invisible. Equation \eqref{eq:tline-half} is the impedance-side statement of the
period $\lambda_g/2$ derived in \eqref{eq:tline-ripple-period}: if the impedance repeats every half wavelength,
then so must the standing-wave pattern, because the pattern is a function of the impedance.
Table~\ref{tab:tline-speciallens} collects the special lengths.

\begin{table}[H]
  \centering
  \caption{Special electrical lengths of a lossless line and the impedance transformation each performs. The periodicity in the last row is the same $\lambda_g/2$ as the standing-wave period of \eqref{eq:tline-ripple-period}.}
  \label{tab:tline-speciallens}
  \small
  \begin{tabular}{@{}lll@{}}
    \toprule
    Electrical length & $\tan\beta d$ & $Z_{\mathrm{in}}$ \\
    \midrule
    $0$                & $0$        & $Z_L$ \\
    $\lambda_g/8$      & $1$        & $Z_0(Z_L+jZ_0)/(Z_0+jZ_L)$ \\
    $\lambda_g/4$      & $\infty$   & $Z_0^2/Z_L$ (short $\leftrightarrow$ open) \\
    $3\lambda_g/8$     & $-1$       & $Z_0(Z_L-jZ_0)/(Z_0-jZ_L)$ \\
    $\lambda_g/2$      & $0$        & $Z_L$ (transparent) \\
    \bottomrule
  \end{tabular}
\end{table}

Both identities are special cases of a single remark: the line is a periodic filter of impedance with period
$\lambda_g/2$, and a quarter wave is the half-period that inverts. In the language of Section~\ref{sec:smith}
\eqref{eq:tline-quarter} is a rotation of the reflection coefficient by $\pi$, which turns $z$ into $1/z$ -- and
a rotation is exactly what a lossless length of line does to $\Gamma$, never a change of $|\Gamma|$. That is
worth stating plainly now, because it is the reason a length of line can never by itself improve a match: it can
move a badly matched impedance to a place where a stub can fix it, but it leaves the distance from the centre of
the chart unchanged.

\subsection{Microstrip as a quasi-TEM line, and the meaning of $\eps_{\mathrm{eff}}$}
\label{sec:tline-epseff}

Everything above assumed a TEM mode: a transverse electric and a transverse magnetic field, both perpendicular
to the direction of propagation, with $E$ and $H$ in a fixed ratio everywhere in the cross-section. A TEM mode
requires a homogeneous dielectric, and this is worth seeing rather than quoting. In a homogeneous medium of
permittivity $\eps$ and permeability $\mu$ a TEM field satisfies $\nabla_t^2 \phi=0$ in the cross-section, the
two-dimensional Laplace equation, and the transverse fields are $-\nabla_t\phi$ and
$-\hat z\times\nabla_t\phi/\eta$ with $\eta=\sqrt{\mu/\eps}$ the wave impedance. A single scalar potential serves
both components, which is why $E/H$ is fixed at $\eta$ everywhere and why $v_p=1/\sqrt{\mu\eps}$ follows without
any assumption about dispersion. Microstrip breaks the homogeneity: the field is partly in the substrate of
permittivity $\eps_r$ and partly in the air above it. The boundary conditions at that interface cannot be met by
a single potential satisfying Laplace's equation in both regions with the same $\beta$, so the true mode is a
hybrid, with a small longitudinal component of both $E$ and $H$.

The engineering response is not to solve the hybrid problem but to define a quantity that makes the TEM
formulae true. This report takes the definition seriously, because the alternative -- writing
$\beta=2\pi f\sqrt{\eps_{\mathrm{eff}}}/c$ as if it were a theorem about microstrip -- is a habit that hides
where the number comes from. Define, for the mode actually propagating,
\begin{equation}
  \eps_{\mathrm{eff}} \equiv \left(\frac{\beta c}{\omega}\right)^2 = \left(\frac{c}{v_p}\right)^2 ,
  \label{eq:tline-epseff}
\end{equation}
with $\beta$ the measured propagation constant and $c$ the speed of light in vacuum. With this definition
\begin{equation}
  \beta = \frac{2\pi f}{c}\sqrt{\eps_{\mathrm{eff}}}
  \label{eq:tline-beta-epseff}
\end{equation}
is an identity, true by construction and carrying no physical content; all the physical content sits in the
claim that a single real number $\eps_{\mathrm{eff}}$, independent of frequency over some band, reproduces the
measured $\beta$ over that band. That claim is an approximation, and Section~\ref{sec:tline-limits} says where it
fails. Operationally \eqref{eq:tline-epseff} says that $\eps_{\mathrm{eff}}$ is not measured directly: it is
extracted from a measured phase, which is exactly what Section~\ref{sec:ms-estimators} does with the phase of
$S_{21}$ across a length of line, and it is why an error in the assumed length becomes an error in
$\eps_{\mathrm{eff}}$.

Given \eqref{eq:tline-epseff}, the characteristic impedance follows from a partial-capacitance argument rather
than from a fit. Fill the cross-section with dielectric of effective permittivity $\eps_{\mathrm{eff}}$ and the
capacitance per unit length becomes $C'=\eps_{\mathrm{eff}}C'_{\mathrm{air}}$, where $C'_{\mathrm{air}}$ is the
capacitance of the same conductor geometry with all dielectric removed. For a non-magnetic substrate
$\mu_r=1$ the inductance per unit length is unchanged by the dielectric, because $L'$ depends on the magnetic
field distribution alone and the dielectric does not perturb it at these frequencies. Substituting into the
lossless identities \eqref{eq:tline-z0-derivation} and \eqref{eq:tline-vp},
\begin{equation}
  Z_0 = \sqrt{\frac{L'}{C'}} = \frac{1}{\sqrt{\eps_{\mathrm{eff}}}}\sqrt{\frac{L'}{C'_{\mathrm{air}}}}
      = \frac{Z_{0,\mathrm{air}}}{\sqrt{\eps_{\mathrm{eff}}}},
  \qquad v_p = \frac{1}{\sqrt{L'C'}} = \frac{c}{\sqrt{\eps_{\mathrm{eff}}}} ,
  \label{eq:tline-z0-epseff}
\end{equation}
which is the reason a microstrip of a given width has a lower impedance on a thicker or higher-permittivity
substrate. The numbers make it concrete: a line of $50\ \Ohm$ with $\eps_{\mathrm{eff}}=3.3$ would be
$50\sqrt{3.3}=90.8\ \Ohm$ in air, so the dielectric does part of the impedance-lowering work and the geometry the
rest. These identities are used without further comment in Section~\ref{sec:ms-z0}, where the same argument is
pushed one step further to the three different impedances that a microstrip measurement can mean, and where the
distinction between them is shown to matter numerically.

\subsection{Where the quasi-TEM picture breaks}
\label{sec:tline-limits}

The quasi-TEM picture is an approximation with identifiable failure modes, and it is worth listing them because
each one shows up somewhere in the measurements of this report.

\emph{Fringing.} The field is not confined to the substrate. The fringing field above the strip makes the
effective width larger than the physical width and makes $C'$ larger than a parallel-plate estimate would give;
it is the physical origin of the fact that $\eps_{\mathrm{eff}}$ lies strictly between $1$ and $\eps_r$ rather
than equalling $\eps_r$. For a narrow strip most of the field is in air and $\eps_{\mathrm{eff}}\to1$; for a
wide strip the field is mostly in the substrate and $\eps_{\mathrm{eff}}\to\eps_r$. Neither limit is a TEM mode,
and the interpolation between them is what an $\eps_{\mathrm{eff}}$ formula approximates.

\emph{Substrate thickness.} If the substrate thickness $h$ becomes comparable to the strip width $w$, the
fringing field is a large fraction of the total and any quasi-static estimate of $C'$ deteriorates rapidly; if
$h$ becomes comparable to $\lambda_g$, the longitudinal field component is no longer small and the hybrid nature
of the mode is no longer a correction.

\emph{Radiation.} An open microstrip line radiates, weakly, from its discontinuities and from any bend; the
radiation is a loss mechanism not captured by $R'$ and $G'$ and it grows with frequency and with substrate
thickness. The important point for this report is the converse, which is a statement about the reference box: a
line inside a closed conducting enclosure cannot radiate into free space at all, because the enclosure's walls
are a boundary condition (the conducting enclosure of the measurement is not the computational domain of
Section~\ref{sec:res-abc}, whose walls are absorbers rather than conductors). Any loss extracted from a
measurement inside a box therefore excludes the radiative channel by construction, and
Section~\ref{sec:ms-quasithem} and Section~\ref{sec:msres-box} discuss what that does to the number that comes
out. It is also why the box is a help and a hindrance at once: it removes a loss mechanism the open line has,
and it adds wall loading the open line does not have.

\emph{Dispersion.} The field distribution shifts between substrate and air as frequency rises -- the mode becomes
slightly more concentrated in the substrate -- so $\eps_{\mathrm{eff}}$ is a weak, increasing function of
frequency. The assumption behind \eqref{eq:tline-beta-epseff}, that one number describes the whole band, is
therefore approximate, and the error grows with $f h$. When $\eps_{\mathrm{eff}}$ is extracted separately at each
frequency and the results drift, that drift is this effect and not necessarily a measurement error; the
broadband extraction of Section~\ref{sec:ms-estimators} is where the distinction has to be made.

\emph{Higher-order modes and conductor thickness.} The lowest substrate mode of a microstrip is not the only one:
the first surface-wave-like mode appears at a frequency set by $h$, and above it the quasi-TEM description is
invalid regardless of how carefully $\eps_{\mathrm{eff}}$ was extracted. Finite conductor thickness, typically
$0.035$ mm of copper, adds a small extra capacitance and a frequency-dependent conductor loss that a
zero-thickness model does not have; at the frequencies of this report it is a correction of a few per cent in
$\alpha$ and a small shift in $\eps_{\mathrm{eff}}$.

The lesson is not that the quasi-TEM picture is wrong but that it is a model with a domain, and the
$\eps_{\mathrm{eff}}$ obtained from a measurement inside a box is the value appropriate to that box. It includes
the wall loading and excludes the radiation, and it is therefore not the open-line value quoted by a synthesis
formula. Reported without that qualification the number would be true and misleading.

\subsection{What the measurement takes from this}
\label{sec:tline-summary}

Three things, and each becomes a chapter. First, the propagation constant: the phase of a travelling wave
advances by $\beta$ per metre, so a pair of voltage samples at known separation yields $\beta$ and hence
$\eps_{\mathrm{eff}}$ through \eqref{eq:tline-epseff} -- which is the whole of the phase-based extraction
discussed in Section~\ref{sec:ms-estimators}, and explains why the ripple of period $\lambda_g/2$ derived in
\eqref{eq:tline-ripple-period} is such a nuisance there. Second, the reflection coefficient:
\eqref{eq:tline-gamma} turns a ratio of two voltages into a single complex number that describes the termination
completely, and the fact that it is bounded by unity for passive loads means it can be drawn inside a disc --
which is the Smith chart of Section~\ref{sec:smith}, and the reason that chart can be read as a map rather than
a table. Third, the characteristic impedance: $Z_0$ is the ratio of voltage to current for a single travelling
wave, and only for a single travelling wave, so any measurement that contains a backward wave must first
separate the two before $Z_0$ means anything at all. That separation is what scattering parameters do, and
Section~\ref{sec:sparams} constructs them.

